\documentclass[11pt]{article}
\usepackage[utf8]{inputenc}
\usepackage[T1]{fontenc}
\usepackage[margin=1.25in]{geometry}
\usepackage{amsmath}
\usepackage{amsfonts}
\usepackage{amssymb}
\usepackage{booktabs}
\usepackage{microtype}
\usepackage{titlesec}
\usepackage{parskip}
\usepackage{abstract}
\usepackage{hyperref}
\usepackage{enumitem}

% Fonts
\usepackage{palatino}
\usepackage{mathpazo}

% Hyperref
\hypersetup{
  colorlinks=true,
  linkcolor=black,
  urlcolor=black,
  citecolor=black
}

% Section formatting
\titleformat{\section}[block]{\large\bfseries\scshape}{\thesection.}{0.75em}{}
\titleformat{\subsection}[block]{\normalsize\bfseries}{\thesubsection.}{0.5em}{}
\titlespacing*{\section}{0pt}{2ex plus 0.5ex minus 0.2ex}{1ex plus 0.2ex}
\titlespacing*{\subsection}{0pt}{1.5ex plus 0.4ex minus 0.2ex}{0.8ex plus 0.1ex}

% Abstract formatting
\renewcommand{\abstractnamefont}{\normalfont\bfseries\scshape}
\setlength{\absleftindent}{1.5em}
\setlength{\absrightindent}{1.5em}

% Custom Commands
\newcommand{\bind}{\operatorname{bind}}
\newcommand{\recomp}{\operatorname{recomp}}
\newcommand{\dsem}{d_{\mathrm{sem}}}
\newcommand{\Jgrowth}{\mathcal{J}_{\mathrm{growth}}}

% Small caps section labels
\setlist[itemize]{leftmargin=1.5em, itemsep=0.25em, topsep=0.5em}

\title{\textsc{The Topology of Intention}\\[0.4em]
{\normalsize\itshape Process-Based Identity and the Geometric Binding of the Self}}
\author{Flyxion}
\date{\today}

\begin{document}

\maketitle

\begin{abstract}
\noindent
What persists when the physical substrate changes? Traditional accounts of identity rely on
material or psychological continuity. This paper develops an alternative: identity as a
\emph{dynamical fixed point}, defined by the persistence of context-sensitive constraint
structures rather than by the conservation of matter or memory. Drawing on Alicia
Juarrero's framework of interlevel causality, we interpret information as the reduction of
possibility, intentions as attractor basins in neural phase space, and the self as a
recursive closure of constraints that reproduces its own organizational conditions under
transformation. The distinction between context-free and context-sensitive constraints is
formalized as a duality between recomposability $\recomp(c)$ and binding force $\bind(c)$,
yielding a geometric picture in which identity occupies regions of deep attractor topology
while noise and decomposability correspond to flat manifolds. We conclude that meaning,
action, and selfhood are unified instances of a single structural phenomenon: the
survival of a constraint across transformation.
\end{abstract}

\bigskip

\section{Introduction: From Force to Constraint}

Modern philosophy of action inherits a tacit commitment to Newtonian efficient causality.
On this picture, intentions act as internal forces: they \emph{push} behaviors from
behind, transferring energy across a causal gap between mind and body. The explanatory
burden falls on specifying the mechanism of impact.

Alicia Juarrero, in \textit{Dynamics in Action} (1999), rejects this framework at the root.
In its place she proposes an account of \textbf{interlevel causality} mediated not by
force but by the operation of \textbf{constraints}. A constraint does not shove a system
from one state to another; it restricts which transitions are possible, thereby reshaping
the system's probability distribution over future states.

This reorientation enables a precise definition of \textbf{information}: not a substance
or content, but the \emph{reduction of possibility} at a source. When the trajectory of a
physical system is fully constrained by the internal organization of a mental structure,
the resulting behavior qualifies as \emph{action}. Action is the unequivocal transmission
of intentional information from a cognitive origin to a behavioral terminus—an information
flow across levels without equivocation.

The present paper develops this framework into a general account of process-based identity.
We formalize the central concepts in the language of dynamical systems, interpret them
geometrically, and show that the persistence of a self is structurally continuous with the
persistence of meaning in discourse and the stability of attractor topology in cognition.
The mechanism is identical across all three domains. Only the substrate differs.

\section{Taxonomy of Constraints}

Juarrero adopts Lila Gatlin's distinction between two constraint types. We render these
formally.

\subsection{Context-Free Constraints}

A \textbf{context-free constraint} moves a system away from equiprobability. Let $\Omega$
be the state space of a system and $P_0$ be the uniform distribution over $\Omega$.
A context-free constraint selects a proper subset $\Omega' \subset \Omega$, reducing
entropy:
\[
H(P_0) > H(P_{\Omega'})
\]
This establishes a baseline of order. The outputs of such a system are \emph{stable} in
the sense that they can be reproduced without reference to context. We associate this
property with \textbf{recomposability}:
\[
\recomp(c) \;\sim\; \text{context-independent stability of output under transformation}
\]
High-$\recomp$ structures are portable and fragment-safe. They survive extraction from
their generating context without significant loss of structure.

\subsection{Context-Sensitive Constraints}

A \textbf{context-sensitive constraint} does more than reduce entropy locally. It
introduces \emph{conditional dependence} across components: the probability distribution
of each part becomes conditioned on the states of other parts. Formally, if $X_i$ are
components of a system, context-sensitive constraints enforce:
\[
P(X_i \mid X_{-i}) \neq P(X_i)
\]
for most $i$, where $X_{-i}$ denotes all components except $i$. This creates a
\emph{global coherence} that is irreducible: decomposing the system into independent parts
destroys the very structure the constraint produces. We associate this with
\textbf{binding force}:
\[
\bind(c) \;\sim\; \text{depth and stability of global constraint coherence}
\]
High-$\bind$ structures resist fragmentation. They are the source of top-down causal
influence: the global organization constrains the behavior of parts in ways that no
local description captures.

\subsection{The Duality}

The two constraint types are not merely different degrees of the same property. They
represent an ontological trade-off:

\begin{center}
\begin{tabular}{@{}lll@{}}
\toprule
\textbf{Property} & \textbf{Context-Free ($\uparrow\recomp$)} & \textbf{Context-Sensitive ($\uparrow\bind$)} \\
\midrule
Stability & Under arbitrary context & Under specific relational configuration \\
Decomposability & Preserved under fragmentation & Lost under fragmentation \\
Information flow & Local and portable & Global and hierarchical \\
Causal structure & Bottom-up & Top-down \\
\bottomrule
\end{tabular}
\end{center}

\noindent
Systems that maximize $\recomp$ sacrifice $\bind$; systems that maximize $\bind$ are
non-decomposable by construction. This trade-off reappears at every level of analysis.

\section{Meaning as Embodied Topology}

Classical information theory, as Shannon developed it, is deliberately semantic-free. A
message is a sequence of symbols; its information content is the reduction in uncertainty
it produces, independent of what it is \emph{about}. This strips the framework of the
capacity to distinguish noise from meaningful signal, or instruction from mere regularity.

Juarrero's dynamical approach reintroduces meaning without mysticism. Meaning, she argues,
is embodied in the \textbf{topographical configuration of phase space}—the arrangement of
attractors, repellers, and basins of attraction that characterize a high-dimensional neural
system.

\subsection{Attractors as Intentions}

An \textbf{attractor} is a region of phase space toward which trajectories converge under
the system's dynamics. A \emph{proximate intention}, in Juarrero's account, functions as
an attractor basin: a valley in the system's mental landscape that draws the trajectories
of motor subsystems into its organizational sphere.

This reframes the question of intentional causation. An intention does not \emph{push}
behavior; it \emph{shapes the landscape} through which behavior moves. The causal relation
is geometric, not mechanical.

\subsection{Binding as Basin Depth}

We can now give $\bind(c)$ a geometric interpretation:
\[
\bind(c) \;\sim\; \text{depth and stability of the attractor basin associated with } c
\]
A highly bound structure occupies a deep attractor. Small perturbations—noise in the
neuromuscular channel, interference from competing intentions—do not dislodge the
trajectory from its constrained path. The intention survives transmission.

Conversely, a low-binding structure corresponds to a shallow attractor: the system is
easily re-routed by contextual fluctuations. The trajectory diverges from its originating
constraint.

Recomposability, in this geometric picture, corresponds to a \emph{flat landscape}: a
manifold without strong attractors where trajectories can be redirected without
resistance. Fragment-safe content inhabits flat territory precisely because it imposes no
strong attractor structure on the reading trajectory.

\subsection{Noise and Equivocation}

Juarrero adopts Shannon's concepts of equivocation and noise to characterize failures of
intentional action. An action fails not because the intention was absent or weak, but
because constraint transmission was disrupted:

\begin{itemize}
  \item \textbf{Noise:} Interference introduced between intention and execution increases
        $\dsem(c, T_\sigma(c))$ and thereby decreases $\bind(c)$. The behavior fails to
        remain bound to its generating constraint.
  \item \textbf{Equivocation:} Ambiguity at the source—an insufficiently organized
        intention—produces multiple competing attractors, allowing the trajectory to
        wander between basins.
\end{itemize}

Intentional action requires that the channel between cognitive source and behavioral
terminus transmit constraint without loss. This is the dynamical analog of high binding.

\section{Identity as a Dynamical Fixed Point}

With this geometric apparatus in place, we can address personal identity directly.

\subsection{The Substance View and Its Limits}

The classical account of personal identity asks what makes person $P$ at time $t_1$ the
same individual as person $P'$ at time $t_2$. Two families of answers dominate:

\begin{itemize}
  \item \textbf{Physical continuity:} $P = P'$ iff they share the same continuous body.
  \item \textbf{Psychological continuity:} $P = P'$ iff $P'$ inherits memories,
        personality, and beliefs from $P$ via the right causal chain.
\end{itemize}

Both views presuppose that identity is a relation between \emph{states}: a snapshot at
$t_1$ is compared with a snapshot at $t_2$, and a criterion is applied. The problem is
that biological and psychological states change continuously. Any fixed threshold for
"enough" continuity is arbitrary.

\subsection{The Process View}

Juarrero proposes that a person is not a thing that changes, but a \textbf{self-organizing
process}—a dynamically stable pattern of information flow, analogous to a whirlpool or a
hurricane. The pattern persists even as its constituent elements are replaced, because
what is conserved is not matter but \emph{organizational structure}.

More precisely: a person is the persistence of a \textbf{context-sensitive constraint
structure} through time. This persistence is not static but dynamic—maintained by the
constant self-reproduction of the constraints that define the system's organization.

\subsection{Recursive Closure}

Juarrero introduces the concept of \textbf{recursive closure} to capture this
self-maintenance. A system achieves recursive closure when its components become mutually
constraining in such a way that the system actively maintains the conditions necessary for
its own existence. The causal structure is circular in a productive sense: the system
produces the constraints that produce the system.

Formally, let $M_t$ denote the constraint structure of the system at time $t$, and let
$\mathrm{dynamics}(\cdot)$ denote the temporal evolution operator. Recursive closure
requires:
\[
M_t \;\approx\; \mathrm{dynamics}(M_t)
\]
This is a \textbf{fixed point condition}. The system is its own attractor. It reproduces
its organizational conditions under the very dynamics those conditions generate.

\subsection{What Persists: The Fixed Point of Constraint}

This fixed point formulation dissolves the substance-based paradox. Identity does not
require that any particular matter, or any particular memory, be preserved. What must be
preserved is the \emph{constraint structure}—the topology of the attractor landscape—that
defines the system's characteristic modes of organization.

We can state this precisely:
\[
\text{Identity} \;=\; \text{constraint structure that satisfies } M_t \approx \mathrm{dynamics}(M_t)
\]
The self is a fixed point of a constraint operator.

\section{Who and What: The Distinction of Trajectories}

\subsection{Two Questions}

Juarrero draws a sharp distinction between two questions that are often conflated:

\begin{itemize}
  \item \textbf{What am I?} — This question concerns the material substrate: the
        biological hardware, the neural tissue, the chemical composition. Answers consist
        of context-free properties, high in $\recomp$ and decomposable without loss.
  \item \textbf{Who am I?} — This question concerns the trajectory: the unique path the
        system has taken through its phase space, the accumulated history of constraints
        that have shaped its current organization.
\end{itemize}

The "What" can be replicated in principle: another arrangement of matter with identical
context-free properties would answer the "What" correctly. The "Who" cannot be replicated,
because it is defined by a \emph{path-dependent history of constraints}. Two systems
cannot share the same trajectory through phase space.

\subsection{Individuality as Path-Dependence}

This provides a mathematical basis for individuality that does not appeal to primitive
haecceity or brute numerical identity. Individuality is path-dependence: no two systems
can have traversed the same trajectory through the space of all possible constraint
configurations. The history is constitutive, not merely biographical.

The growth functional $\Jgrowth$ tracks the construction of stable constraint topology
over time:
\[
\Jgrowth \;\sim\; \text{construction of stable attractor topology through recursive constraint formation}
\]
A life, in this formalism, is a trajectory through constraint space that progressively
deepens the basin structure of the system's organization—increasing $\bind(M_t)$ through
the accumulated residue of experience.

\section{Relational Identity and Plasticity}

\subsection{Open Systems}

Juarrero's account has an important consequence for the spatial locus of identity. A
self-organizing system is necessarily an \emph{open system}: it exchanges matter, energy,
and information with its environment. The constraints that constitute the self are not
sealed inside the skull; they are defined partly by the system's characteristic modes of
interaction with its environment.

Identity, on this view, is \textbf{relational}: it is not a property of the organism in
isolation but of the organism-environment coupling. The constraint structure that defines
"Who I am" includes stable patterns of environmental engagement—habitual responses,
skilled practices, long-standing relationships.

\subsection{Plasticity and Bifurcation}

This relational picture accommodates radical change without identity loss. A complex
dynamical system can undergo a \textbf{bifurcation}: a qualitative reorganization of its
attractor structure in response to a sufficiently strong perturbation. After a bifurcation,
the system occupies a different region of phase space—a different basin of attraction.

Yet the system that undergoes the bifurcation is the same system, because the bifurcation
is a transformation of a continuous trajectory, not a discontinuous substitution of one
system for another. The path-dependence is preserved across the bifurcation event.

This is the dynamical analog of significant personal transformation: the person who
emerges from a major life change is not the same in every respect as the person who
entered it, but the transformation is a reorganization of the same constraint history,
not its erasure.

\section{The Unified Picture}

We are now in a position to state the central thesis in its most general form.

\subsection{Three Domains, One Mechanism}

At three levels of analysis—discourse, cognition, and selfhood—the same structural
phenomenon governs the persistence of organization:

\begin{enumerate}
  \item \textbf{Discourse:} Semantic content with high $\bind(c)$ survives transformation
        (extraction, recombination, adversarial reading) without loss of meaning. Content
        with high $\recomp(c)$ survives fragmentation but loses coherence when the
        context-sensitive structure is disrupted.

  \item \textbf{Cognition:} Intentional action requires that the constraint structure of
        the intention survive transmission through the neuromuscular channel without
        equivocation. The intention is an attractor; its depth determines whether the
        trajectory remains bound.

  \item \textbf{Selfhood:} Personal identity requires that the constraint structure of
        the self satisfy a fixed point condition under its own dynamics. The self is a
        recursive closure: a process that maintains the conditions of its own organization.
\end{enumerate}

The mechanism is identical. The constraint is the cause; what varies is only the level
at which the constraint operates.

\subsection{A Unified Principle}

\begin{quote}
\itshape
Any system governed by constraints evolves by reducing its space of possible trajectories.
The persistence of structure—whether semantic, behavioral, or personal—depends on the
ability of those constraints to survive transformation across levels. The difference
between shallow and deep systems, between noise and meaning, between behavior and action,
is not a matter of content. It is a matter of how strongly constraints bind trajectories
across transformation.
\end{quote}

\section{Conclusion}

We began with a question about causality and arrived at an account of the self. This
trajectory is not accidental. The two questions are structurally connected: what it means
for an intention to cause an action is the same, at a different scale, as what it means
for a self to persist through time.

In both cases, the answer invokes the same formal structure: a constraint that survives
transformation because the system has achieved sufficient recursive closure to reproduce
the conditions of its own organization.

Juarrero's contribution is to show that this structure is not mysterious—it is the
ordinary operation of dynamical systems with sufficiently rich attractor topology. Meaning
is not added to the world from outside; it is embodied in the shape of the landscape
through which trajectories move. The self is not a ghost in a machine; it is a fixed
point of the machine's own constraint dynamics.

What is conserved across change is not matter, not memory, but the capacity to regenerate
the same organizational conditions from within. This is what it is to be, in Juarrero's
precise sense, a \emph{self}:

\[
\boxed{M_t \;\approx\; \mathrm{dynamics}(M_t)}
\]

The constraint is the cause. The trajectory is the identity. The fixed point is the self.

\end{document}
